% Offizielle Beispieldatei für beamer-Vorlage aus tubslatex Version 0.3beta2
\documentclass[fleqn,11pt,aspectratio=43]{beamer}

\usepackage[ngerman]{babel}
\usepackage[utf8x]{inputenc}
\usepackage{graphicx}
\usepackage{listings}
\usepackage{color}
\usepackage{pxfonts}

\usetheme[%
  %cmyk,%<cmyk/rgbprint>,          Auswahl des Farbmodells
  orange,%<blue/orange/green/violet> Auswahl des Sekundärfarbklangs
  %dark,%<light,medium>        Auswahl der Helligkeit
  %colorhead,%    Farbig hinterlegte Kopfleiste
  %colorfoot,%    Farbig hinterlegt Fußleiste auf Titelseite
  colorblocks,%   Blöcke Farbig hinterlegen
  %nopagenum,%    Keine Seitennumer in Fußzeile
  %nodate,%       Kein Datum in Fußleiste
  %tocinheader,%   Inhaltsverzeichnis in Kopfleiste
  %tinytocinheader,% kleines Kopfleisten-Inhaltsverzeichnis
  %widetoc,%      breites Kopfleisten-Inhaltsverzeichnis
  %narrowtoc,%    schmales Kopfleisten-Inhaltsverzeichnis
  %nosubsectionsinheader,%  Keine subsections im Kopfleisten-Inhaltsverzeichnis
  %nologoinfoot,% Kein Logo im Fußbereich darstellen
  ]{tubs}

% Titelgrafik, automatisch beschnitten, Weitere Optionen: <scaled/cropx/cropy>
%\titlegraphic[scaled]{\includegraphics{img/titlepicture.jpg}}
\titlegraphic[cropy]{\includegraphics{img/zynq.jpg}}

% Logo, dass auf Titelseiten oben rechts und auf Inthaltsseiten unten rechts
% dargestellt wird. Es wird jeweils automatisch skliert
\logo{\includegraphics{img/logo_mit_text.pdf}}

% settings for source code highlighting
\lstset{ %
  backgroundcolor=\color{white},         % choose the background color
  basicstyle=\ttfamily\footnotesize,     % size of fonts used for the code
  breaklines=true,                       % automatic line breaking only at whitespace
  captionpos=b,                          % sets the caption-position to bottom
  commentstyle=\color{tuGreenDark80},    % comment style
  escapeinside={\%*}{*)},                % if you want to add LaTeX within your code
  keywordstyle=\color{tuBlueMedium100}\bfseries, % keyword style
  stringstyle=\color{tuVioletMedium},    % string literal style
}

% Titelseite
\title{SoCRocket Week 2014}
\subtitle{Introduction of SoCRocket}
\author{Rolf Meyer}

\begin{document}

\begin{frame}[plain]
\titlepage
\end{frame}

\part{Introduction}

\begin{frame}[plain]
  \partpage
\end{frame}
\begin{frame}{Contens}
\tableofcontents
\end{frame}

\section{File structure}

\begin{frame}{File structure - 1}

\begin{itemize}
  \item /contrib
  \begin{itemize}
    \item Patches for external dependencies
    \item Rewritten dependencies
    \item Own dependencies
  \end{itemize}
  \item /common
  \begin{itemize}
    \item Library for common used classes and functions
    \item Small platform helpers
    \item VerbosityKit, MsgLogger, Power, EndianessHelper
  \end{itemize}
  \item /models
  \begin{itemize}
    \item Base directory for all models
    \item Each model is a library in an own subdirectory
    \item AHBCtrl, APBCtrl, AHBMem, Irqmp, ...
  \end{itemize}
\end{itemize}
\end{frame}

\begin{frame}{File structure - 2}

\begin{itemize}
  \item /models/utils
  \begin{itemize}
    \item Suporting infrastructure to build models
    \item APBDevice, AHBDevice, AHBMaster, AHBSlave, MEMDevice
  \end{itemize}
  \item /platforms
  \begin{itemize}
    \item Basedirectory for all paltforms
    \item All platforms formed from the models are stored in this directory
    \item leon3mp.platform, yuvstream.platform, mips1sp.platform
  \end{itemize}
  \item /templates
  \begin{itemize}
    \item Basedirectory for the configuration files of a platform
    \item Needed for the generator wizard
    \item XML Templates (.tpa), Configfile (.json)
  \end{itemize}
\end{itemize}
\end{frame}

\begin{frame}{File structure - 3}

\begin{itemize}
  \item /adapters
  \begin{itemize}
    \item VHDL-Transactors
  \end{itemize}
  \item /tools
  \begin{itemize}
    \item Home to all scripts
    \item Including generator wizard and build system scripts
    \item get\_json\_attr, set\_json\_attr, list.sh
  \end{itemize}
\end{itemize}
\end{frame}

\section{Python - Crash course}

\begin{frame}{Python - Explaination of the Syntax}
\begin{itemize}
  \item Most languages use curly braces to group pieces of code. 
  \item Like everything in a for loop will be surrounded by a \{ and a \}. 
  \item Not in Python. 
  \item Indention is the key. 
  \item Everything on the same indention level is part of the same group of code.
\end{itemize}
\end{frame}

\begin{frame}[fragile]{Python - Loops}
\begin{lstlisting}[language=python]
for loop stuff:
    weee
    fun
this_isnt_part_of_the_loop()
\end{lstlisting}
Each statement is (usually) on its own line whereas a semicolon is used to delimit lines in most other programming language syntaxes.
\end{frame}

\begin{frame}[fragile]{Python - Functions}
Functions are defined using the def statment. It goes like this...
\begin{lstlisting}[language=python]
def functionname(parameter1, parameter2):
    code goes here
    and here 
    etc.
\end{lstlisting}
\end{frame}

\begin{frame}[fragile]{Python - Classes}
Classes are defined by the class statement.
\begin{lstlisting}[language=python]
class Classname(ParentClass):
    def __init__(self):
        self.member = 5
\end{lstlisting}
Member functions always have the first prameter self. It is the equivalent to this.
\end{frame}

\begin{frame}[fragile]{Python - Beware the Indetion!}
And you can have more indentions under a def in the same way you would have used curly braces before...
\begin{lstlisting}[language=python]
def foo():
    do stuff
    do more stuff
    for loop time:
        this is part of the loop
        so is this
    but this is not, this is just part of the rest of the function
\end{lstlisting}
\end{frame}

\begin{frame}[fragile]{Python - Lists}
\begin{itemize}
  \item Lists are very versatile.
  \item Lists do not all have to contain the same type of data. 
  \item It's ok to mix strings and numbers and objects and nested lists, etc.
\end{itemize}
\begin{lstlisting}[language=python]
shopping_list = ["bread","milk","hamster brains"]
\end{lstlisting}

then if you want to add an item to the list...
\begin{lstlisting}[language=python]
shopping_list = shopping_list + ["turnips"]
\end{lstlisting}

\end{frame}

\begin{frame}[fragile]{Python - Lists access}
If you want to access just one item in that list you can do it like this...
\begin{lstlisting}[language=python]
shopping_list[0]
\end{lstlisting}
That's the first item in the list. 
\begin{lstlisting}[language=python]
shopping_list[1]
\end{lstlisting}
would be the 2nd, etc.
\end{frame}

\begin{frame}[fragile]{Python - Lists range}
If you want a range of items in the list you can do that with colon magic...
\begin{lstlisting}[language=python]
shopping_list[0:3]
\end{lstlisting}
Accesses the first 3 items in the list. 
\begin{itemize}
  \item The way it works is the first number is where you want to start
  \item The second number is where you want to end (not how many you want).
  \item It goes up to but not including the 2nd number.
  \item Leave them empty to go frome the start/to the end.
  \item {\tt shopping\_list[:]} gets a copy of the hole list.
\end{itemize}
\end{frame}

\begin{frame}[fragile]{Python - For loops}
This is the biggest change from most languages:
\begin{itemize}
  \item Traditionally, a variable increments until a condition is false. 
\end{itemize}
For loops in Python:
\begin{itemize}
  \item The interpreter goes throug to each item of a list with a variable
  \item Making the variable equal to each one as it goes...
\end{itemize}
\begin{lstlisting}[language=python]
for thingy in shoppin_list:
    buy(thingy)
\end{lstlisting}

there's a loop (assuming we're using the same shopping list from earlier)
\end{frame}

\begin{frame}[fragile]{Python - Loop over integers}
If you want a list of numbers that go from x to y, here's how you do that:
\begin{lstlisting}[language=python]
list_of_nums = range(x, y + 1)
\end{lstlisting}

\begin{itemize}
  \item Then list of nums will be a list of numbers...in that range. 
  \item It always starts at the first number and goes up to (but not including) the 2nd number. 
  \item That's what that + 1 thing is in there for.
\end{itemize}
\end{frame}

\begin{frame}[fragile]{Python - If statements}
\begin{lstlisting}[language=python]
if some condition:
    do this
\end{lstlisting}

Simple enough:

\begin{lstlisting}[language=python]
if x > 4:
    print "OMG, x must be HUGE!"
\end{lstlisting}
\end{frame}

\begin{frame}[fragile]{Python - Or else!}
Then there's else:
\begin{lstlisting}[language=python]
if x > 4:
    print "OMG, x must be HUGE!"
else:
    print "x ain't nuttin'"
\end{lstlisting}
\end{frame}

\begin{frame}[fragile]{Python - Or else if!}
else if's are denoted by elif's

\begin{lstlisting}[language=python]
if x > 4:
    print "OMG, x must be HUGE!"
elif x == 4:
    print "x is four"
else:
    print "x is tiny"
\end{lstlisting}
\end{frame}


\begin{frame}[fragile]{Python - While loops}
While loops look just like if's and work like how you would imagine them
\begin{lstlisting}[language=python]
while condition:
    do this
\end{lstlisting}
\end{frame}


\begin{frame}[fragile]{Python - Strings}
Strings are concatenated with + signs
\begin{lstlisting}[language=python]
greeting = "Hello, " + "World!"
\end{lstlisting}
To define a string you can use ", ', """ or '''. This saves the hassle of escaping.
And vprintf replacment ist build in ({\tt \%s, \%d, \%f...}):
\begin{lstlisting}[language=python]
name = "%(first)s %(last)s" %  {
    'first':name[0], 
    'last':name[1]
}
greeting = """Hello %s""" % name
\end{lstlisting}
\end{frame}

\begin{frame}[fragile]{Python - == and =}
== and = are just like most languages:
\begin{lstlisting}[language=python]
foo = "yo"
if (foo == "yo"):
    print "foo is yo"
\end{lstlisting}
\end{frame}


\begin{frame}[fragile]{Python - Comments}
Comments are done with a \#:
\begin{lstlisting}[language=python]
print "this is some code" #this line prints stuff
\end{lstlisting}
\end{frame}


\begin{frame}[fragile]{Python - And and Or}
  There are no \&\& or ||, it's just {\tt and} and {\tt or}
\begin{lstlisting}[language=python]
if (x == 2) and (y > 5):
    do_stuff()
\end{lstlisting}
\end{frame}

\begin{frame}[fragile]{Python - But there is much more}
  \begin{itemize}
    \item There are tonns of usefull features in the python language.
    \begin{itemize}
      \item Go to "https://docs.python.org/2/tutorial/index.html" for a full tutorial.
      \item Language references you can finde here:
      \item https://docs.python.org/2/reference/index.html
    \end{itemize}
    \item And like Java Python comes with a standard libray
    \begin{itemize}
      \item Which is full of useful things
      \item os, json, base64, HTMLParser, shlex, etc.
      \item https://docs.python.org/2/library/index.html
    \end{itemize}
  \end{itemize}
\end{frame}

\section{Build System}

\begin{frame}[fragile]{Build System}
  \begin{itemize}
    \item Waf 1.7
    \item Python based
    \item Split into serveral files: “wscript”s and python scripts!
    \item It’s simple: 
  \end{itemize}
\begin{lstlisting}[language=bash]
./configure; make; make install
\end{lstlisting}
For waf: 
\begin{lstlisting}[language=bash]
./waf configure; ./waf build; ./waf install
./waf configure build install # in short
\end{lstlisting}
\end{frame}

\begin{frame}[fragile]{Build System}
  \begin{itemize}
    \item Build ist default!
    \item To reduce build time compile only needed targets
  \end{itemize}
  \begin{lstlisting}[language=bash]
./waf --target=<target>
  \end{lstlisting}

Usefull:
  \begin{lstlisting}[language=bash]
./waf list      # List all targets
./waf clean     # Clean build dir
./waf distclean # Remove build dir
  \end{lstlisting}
\end{frame}

\subsection{Waf}

\begin{frame}[fragile]{Waf - Top-level wscript}
The top-level wscript is the collection point to a common build system and initializes all the tools needed to build.

Most of them are outsourced in python scripts which residents in /tools/waf/
The only remains are the phase functions and call hooks:

\begin{lstlisting}[language=python]
#! /usr/bin/env python
# vim : set fileencoding=utf-8 expandtab noai ts=4 sw=4 filetype=python :
APPNAME = 'SoCRocket'
top = '.'
out = 'build'

import sys
from waflib.Tools import waf_unit_test
from tools.waf.logger import Logger

LOAD = ['compiler_c', 'compiler_cxx', 'python', 
        'swig', 'waf_unit_test', ...]
\end{lstlisting}
\end{frame}

\begin{frame}[fragile]{Waf - Top-level wscript 2}
\begin{lstlisting}[language=python]
TOOLS = ['common', 'flags', 'virtualenv', 'pthreads', 
         'boost', 'modelsim', 'systools', ...]

def options(self):
    self.load(LOAD)
    self.load(TOOLS, tooldir='tools/waf')

def configure(self):
    self.load(LOAD)
    self.check_waf_version(mini='1.6.0')
    self.check_python_version((2,4,0))
    self.check_python_headers()
    self.load(TOOLS, tooldir='tools/waf')

def build(self):
    self.recurse_all()
\end{lstlisting}
\end{frame}

\begin{frame}[fragile]{Waf - Recursion wscripts}
Some wscripts are just needed for reaching further subdirectories:
Models, platforms, software, …
  
\begin{lstlisting}[language=python]
#! /usr/bin/env python
# vim : set fileencoding=utf-8 expandtab noai ts=4 sw=4 filetype=python :
top = '..'

def build(self):
  self.recurse_all()
  self.recurse_all_tests()
\end{lstlisting}
The "recurese\_*"-members come from the common file (/tools/waf) included in the top-level wscript as first tool
\end{frame}

\begin{frame}[fragile]{Waf - Models wscript}
Each models is a static C++ Libary! Nothing more  
  
\begin{lstlisting}[language=python]
#! /usr/bin/env python
# vim : set fileencoding=utf-8 expandtab noai ts=4 sw=4 filetype=python :
top = '../..'

def build(self):
  bld(
    target          = 'gptimer',
    features        = 'cxx cxxstlib',
    source          = 'gptimer.cpp gpcounter.cpp',
    export_includes = '.',
    uselib          = 'BOOST SYSTEMC TLM AMBA GREENSOCS',
    use             = 'common signalkit utils',
    install_path   = '${PREFIX}/lib',
  )
\end{lstlisting}
\end{frame}

\begin{frame}[fragile]{Waf - Platforms wscript}
Tests or Platforms are C++ applications linked with all libraries needed for execution
  
\begin{lstlisting}[language=python]
#! /usr/bin/env python
# vim : set fileencoding=utf-8 expandtab noai ts=4 sw=4 filetype=python :
top = '../..'

def build(self):
    bld(
        target       = 'leon3mp.platform',
        features     = 'cxx cprogram pyembed',
        source       = 'sc_main.cpp',
        includes     = '.',
        use          = 'ahbctrl irqmp gptimer apbctrl ...'
        uselib       = 'TRAP BOOST ELF_LIB SYSTEMC ...'
    )
\end{lstlisting}
\end{frame}

\begin{frame}[fragile]{Waf – Target software wscripts}
Waf can be used to compile Target software exactly the same as for host software (Models/Libraries or Platforms/Applications):
  
\begin{lstlisting}[language=python]
#! /usr/bin/env python
# vim : set fileencoding=utf-8 expandtab noai ts=4 sw=4 filetype=python :
top = '../..'
def build(self):
  bld(
     features     = 'c cprogram sparc',
     target       = 'hello.sparc',
     cflags       = '-static -g -O1 -lm -mno-fpu',
     linkflags    = '-static -g -O1 -lm -mno-fpu',
     source       = ['hello.c'],
     install_path = None,
  )
\end{lstlisting}
The only difference is the feature {\tt sparc}
\end{frame}

\begin{frame}[fragile]{Waf – Unittests}
If a test shall be performed on a host target just add the feature {\tt test}.
If it needs to be executed in a specific platform (host and target code) use a systest:
\begin{lstlisting}[language=python]
#! /usr/bin/env python
# vim : set fileencoding=utf-8 expandtab noai ts=4 sw=4 filetype=python :
top = '../..'

def build(self):
  bld(
      features    = 'systest',
      system      = 'leon3mp.platform',
      rom         = 'sdram.prom',
      ram         = 'hello.sparc',
  )
\end{lstlisting}
\end{frame}

\begin{frame}[fragile]{Waf – Modelsim Targets}
The feature modelsim uses modelsim to compile. You can specify it as a unittest too.
\begin{lstlisting}[language=python]
#! /usr/bin/env python
# vim : set fileencoding=utf-8 expandtab noai ts=4 sw=4 ...
top = '..'
def build(self):
  bld(
    target   = 'gptimer.1.rtl.test',
    features = 'modelsim test',
    source   = ['gptimer_wrapper.vhd',
                'test1/testbench.cpp', 'test1/top_rtl.cpp',
                '../../../adapters/rtltransactor.cpp',
                '../../utils/clkdevice.cpp',
                '../../utils/apbdevice.cpp'],
    uselib   = 'AMBA GREENSOCS BOOST',
    includes = '. test1 ../../utils',
  )
\end{lstlisting}
\end{frame}


\subsection{Python tools}
\begin{frame}[fragile]{Waf Python Tools - Common}
Common registes a basic set of tools to compile SoCRocket.
\begin{itemize}
  \item The recursion utilities
  \item helper for fetching and compiling dependencies
\end{itemize}
\end{frame}

\begin{frame}[fragile]{Waf Python Tools - Configuration}
In the configuration phase all dependencies are checked and stored 
\begin{lstlisting}[language=python]
#! /usr/bin/env python
# vim : set fileencoding=utf-8 expandtab noai ts=4 sw=4 filetype=python :
top = '..'

def find(self):
    self.check_cxx(
        lib='elf',
        uselib_store='ELF_LIB',
        mandatory=True,
        libpath = libpath,
        errmsg  = "not found"
    )

def configure(self):
    find(self)
\end{lstlisting}
\end{frame}

\begin{frame}[fragile]{Waf Python Tools - Deps}
Moreover socrocket is fetiching missing dependencies (via tar or git):
\begin{lstlisting}[language=python]
#! /usr/bin/env python
# vim : set fileencoding=utf-8 expandtab noai ts=4 sw=4 ...
top = '..'

def configure(self):
    try:
        find(self)
    except:
        name    = "libelf"
        version = "0.8.13"
        self.dep_build(
            name    = name, version = version,
            tar_url = "http://www.mr511.de/%(base)s.tar.gz",
        )
        find(self, self.dep_path(name, version))
\end{lstlisting}
\end{frame}

\part{ToDos}

\begin{frame}[plain]
  \partpage
\end{frame}
\begin{frame}{Contens}
\tableofcontents
\end{frame}

\section{What needs to be done}

\begin{frame}{ToDos}
  \begin{itemize}
    \item Lorem ipsum dolor sit amet, consetetur sadipscing elitr, sed diam
      nonumy eirmod tempor invidunt ut labore et dolore magna aliquyam
    \item At vero eos et accusam et justo duo dolores et ea rebum.
      \begin{itemize}
        \item Stet clita kasd gubergren, no sea takimata sanctus est Lorem ipsum
          dolor sit amet!
          \begin{itemize}
            \item Nam eget dui.
            \item Maecenas tempus, tellus eget condimentum rhoncus, sem quam
              semper libero, sit amet adipiscing sem neque sed ipsum.
          \end{itemize}
        \item Duis leo
      \end{itemize}
    \item Aliquam lorem ante, dapibus in, viverra quis, feugiat a, tellus. 
  \end{itemize}
\end{frame}

\begin{frame}{SoCRocketDevice}
  \begin{itemize}
    \item A common parant device to derive from.
    \item All models should derive from it (over AHBDevice, APBDevice)
    \item A common class to implement API calls
  \end{itemize}
\end{frame}

\begin{frame}{GS\_Config Parameter}
  \begin{itemize}
    \item All platforms use gs\_params to implement component generics like in grlib designs.
    \item But these parameters are generics of the components/models. 
    \item They should be implemented inside the models.
    \item A initializer function called {\tt init\_configuration()} should be present in each model.
    \item In this method all descriptions and attributes from the {\tt .tpa}-file should be added.
  \end{itemize}
\end{frame}

\begin{frame}{Register Init}
  \begin{itemize}
    \item Like the gs\_config specific function {\tt init\_configuration()} a 
    \item In this method all descriptions and attributes from the {\tt .tpa}-file should be added.
  \end{itemize}
\end{frame}

\begin{frame}{sc\_report VerbosityKit backend }
\end{frame}

\begin{frame}{Documentation}
\begin{itemize}
    \item Update Doxygen documentation from the user manual.
    \item Each directory needs to have a {\tt <model>.md}
    \item {\tt doc/*.md} for common chapters from the user manual and hldse
\end{itemize}
\end{frame}

\begin{frame}{Documentation guid}
\end{frame}

\begin{frame}{Coding Style}
We will follow the google coding guidlines
(It's not perfect but already existing, including tools)
\begin{itemize}
  \item http://google-styleguide.googlecode.com/svn/trunk/cppguide.xml
\end{itemize}
Make use of coding style tools like:
\begin{itemize}
  \item Cpplint.py (from goole)
  \item OCLint
  \item clang sanitizer
\end{itemize}
\end{frame}

\begin{frame}[fragile]{\#define Guards}
\begin{itemize}
  \item All header files should have \#define guards to prevent multiple inclusion.
  \item The format of the symbol name should be {\tt <PROJECT>\_<PATH>\_<FILE>\_H\_}.
  \item To guarantee uniqueness, they should be based on the full path in a project's source tree. 
  \item For example, the file foo/src/bar/baz.h in project foo should have the following guard:
\end{itemize}

\begin{lstlisting}[language=c]
#ifndef SOCROCKET_MODELS_GPTIMER_GPTIMER_H_
#define SOCROCKET_MODELS_GPTIMER_GPTIMER_H_
//...
#endif  // SOCROCKET_MODELS_GPTIMER_GPTIMER_H_
\end{lstlisting}
\end{frame}

\begin{frame}[fragile]{includes}
  \begin{itemize}
    \item All of a project's header files should be listed as descendants of the project's source directory.
    \item Without use of UNIX directory shortcuts . (the current directory) or .. (the parent directory). 
  \end{itemize}

For example, socrocket/models/gptimer/gptimer.h should be included as

\begin{lstlisting}[language=c]
#include "core/models/gptimer/gptimer.h"
\end{lstlisting}

  \begin{alertblock}{This will break our project structure}
    But everythng will get more consistent. \\
    every target needs to include self.root (waf)
  \end{alertblock}
\end{frame}

\end{document}
